\documentclass[10pt,a4paper]{amsart}
\usepackage[margin=19mm]{geometry}
\usepackage{amsmath,amssymb,mathtools}
\usepackage[scr=rsfs,cal=euler]{mathalfa}
\usepackage{xfrac}
\usepackage[unicode,hidelinks,bookmarks=false]{hyperref}
\newcommand{\ie}{i.e.\ }
\newcommand{\cf}{cf.\ }
\newcommand{\resp}{respectively\ }
\newcommand{\Subsection}{\subsection}
\providecommand{\nobreakdash}{\nobreak\hbox{-}}
\setlength{\emergencystretch}{2em}
\multlinegap0pt
\usepackage[normalem]{ulem}
\usepackage{amssymb}
\usepackage{xcolor}
\usepackage[shortlabels]{enumitem}
\usepackage{bm}
\usepackage{mathtools}
\usepackage{nicefrac}
\usepackage{algorithm}\usepackage{algpseudocode}
\usepackage{float}
\newtheorem{lemma}{Lemma}
\newtheorem{theorem}{Theorem}
\newcommand{\LG}{\check{G}}
\newcommand{\UG}{\hat{G}}
\newcommand{\nicetrees}{T}
\title{An improvement on the bound for the acyclic chromatic index}
\author{Lefteris Kirousis, John Livieratos,, Alexandros Singh}
\date{Source: arXiv:2602.14859v1 (CC BY 4.0)}
\begin{document}
\maketitle

\noindent\emph{Source and licence.} An improvement on the bound for the acyclic chromatic index. Version arXiv:2602.14859v1 (CC BY 4.0), distributed by arXiv under the Creative Commons Attribution 4.0 International licence, http://creativecommons.org/licenses/by/4.0/, as recorded in the arXiv licence metadata for this exact version and shown on the abstract page https://arxiv.org/abs/2602.14859v1. Full licence notice: \texttt{licenses/arxiv-2602.14859-CC-BY-4.0.txt}.\par\smallskip

\noindent\textit{Editorial scope.} Counted unit: the article's theorem thm:boundstoradius (source0.tex lines 957--962). The article's complete original proof of it is delivered with it (source0.tex lines 963--1071, one \texttt{proof} environment of 109 lines). No mathematics is written, repaired or re-derived: the statement and the proof are the article's own lines, in the article's own notation, and no retained line is substituted or re-flowed.  Retained uncounted as the source-local context the proof needs: lines 774--780; lines 841--843; lines 906--910; lines 918--930.  Nothing was omitted from any retained span, so the package carries no omission record.  No retained line contains a citation, so the wrapper carries no bibliography and no bibliography entry.  No retained line contains a figure environment, an image-inclusion command or drawing code, so no image is delivered and the package declares no asset dependency.  Editorial formatting only, in addition: an AMS \texttt{amsart} preamble that restates the article's own declarations and macros used by the retained lines, namely the environments \texttt{lemma}, \texttt{theorem} on separate counters, together with the standard packages those lines need.  The retained environments print in this document as Lemma 1, Theorem 1 in order of appearance, whereas the article prints the same items with the numbering of its own full document.  The complete native source of the article as distributed in the arXiv e-print for this exact version is bundled unchanged as \texttt{sources/194751/source0.tex}.
\begin{lemma}\label{lemma:spec1}
The generating function $\nicetrees(z)$ obeys the functional equation:
\begin{equation}\label{eq:spec}
 T(z) = \frac{z}{2} \exp\left(\sum\limits_{j \geq 1} \frac{\nicetrees(z^j)}{j}\right) + \frac{z}{2}\exp\left(\sum\limits_{j \geq 1} (-1)^j \frac{\nicetrees(z^j)}{j} \right)  
-\frac{z}{2}\left({T^2(z)} + {\nicetrees(z^2)} \right).
\end{equation}
\end{lemma}

\begin{lemma}\label{lem:planeTreeCoeffIneq} $B(z)$ dominates coefficient-wise C(z), i.e.,
 $[z^n] C(z) \leq [z^n] B(z)$ and $C(z)$ dominates coefficient-wise $T(z).$
\end{lemma}

\begin{lemma}\label{lemma:powerbounds}
Let $f(z)$ be a power series with nonnegative coefficients and
$[z^0] f(z) = 0$, $[z^1] f(z) = 1$ and radius of convergence $0 < r < 1$. 
Then $f(x^i) \leq f(x^2)^{i/2}$ for $i \geq 2, x \in [0,r^{1/2})$.
\end{lemma}

\begin{lemma}\label{lemma:implicitsing}
Let $G(z,w)$ be an entire  function such that:
\begin{itemize}
    \item $[z^nw^m]G(z,w) \geq 0$ and nonzero for some $m\geq 1$,
    \item $G(0,w) = 0$, for all $w$,
 \end{itemize}
  Then there exists an analytic solution $y(z) = \sum_{n=1}^{\infty}a_nz^n$ to the functional
equation $y=G(z,y)$ with $y(0)=0$, with uniquely determined $a_n \geq 0$,   and  positive finite radius of convergence $r$, such that   $\lim_{x \rightarrow r^{-}} y(x) = s$ exists and is finite.  Moreover $r,s$ satisfy the {\em characteristic} system,
\begin{equation*}
    G(r,s) = s,  \ 
    G_w(r,s) = 1.
\end{equation*}
\end{lemma}

\begin{theorem}\label{thm:boundstoradius}
The radius convergence $\rho$ of  $T(z)$ satisfies: 
\begin{equation*}
    0.6677 \leq \rho \leq 0.6678.
\end{equation*}
\end{theorem}

\begin{proof}


At first, consider the two variable function:

\begin{align}\label{eq:G}
\nonumber G(z,w) =   &\frac{z}{2} 
\ \exp\left( w + \sum\limits_{j \geq 2} \frac{\nicetrees(z^j)}{j}\right) 
+ \frac{z}{2}\exp\left(-w + \sum\limits_{j \geq 2} (-1)^j \frac{\nicetrees(z^j)}{j}\right) \\ &
- \frac{z}{2}\left({w^2}
+ {\nicetrees(z^2)} \right),
\end{align}
and we observe that by Lemma \ref{lemma:spec1}, $T(z) = G(z,T(z))$.

Let $N$ be a fixed (large) even integer.

We begin with computing the upper bound.   Recalling that we can compute coefficients of $T(z)$ up to any order, we compute a truncation of $T(z)$ to powers at most  $N$, which we call $L(z)$ ($L(z)$ is a polynomial of degree $N-1$, since even powers of $z$ in $T(z)$ vanish ---there exist no rooted trees with even outdegrees and even number of nodes). Then in $G$ above, we plug  $L(z^j)$ for $T(z^j)$  and restrict the sums in the expoentials to $j=2, \ldots, N$. We thus obtain some  $\UG(z,w)$, where:
\begin{align}\label{eq:UG}
\nonumber \UG(z,w) =   &\frac{z}{2} 
\ \exp\left( w + \sum\limits_{j =2}^{N} \frac{L(z^j)}{j}\right) 
+ \frac{z}{2}\exp\left(-w + \sum\limits_{j =2}^{N} (-1)^j \frac{L(z^j)}{j}\right) \\ &
- \frac{z}{2}\left({w^2}
+ {L(z^2)} \right).
\end{align}
We claim that $\UG(z,w)$ satisfies  all conditions for Lemma \ref{lemma:implicitsing}.  Indeed, first  observe that 
$\UG(z,w)$  is  entire,  because of the truncations. The first bullet in Lemma~\ref{lemma:implicitsing} follows by expanding the two exponentials in $\UG(z,w)$ as power series and observing, as before,  that terms with negative signs either cancel out or are absorbed  by ones with positive signs. The second  bullet  is obvious.  So, let $r,s$ be the positive real numbers that satisfy 
$$\UG(r,s)=r,\UG_w(r,s)=1.$$
 Observe now that  the solution $\hat{y}(z)$ to $\hat{y}(z) = \UG(z,\hat{y}(z))$
satisfies $[z^n] \hat{y}(z) \leq [z^n] \nicetrees(z)$, because  $\UG$ was obtained by truncations in $G$ and both  $[z^n] \hat{y}(z)$ and  $[z^n] \nicetrees(z), n\geq 1,$ are  obtained by polynomial recurrences with nonnegative coefficients defined via $\hat{y} = \UG(z,\hat{y}(z))$ and $T(z) = G(z, T(z))$, respectively. So $r$ is an upper bound on $\rho$. Finally, we choose, for example,  $N=100$, and  use a computer algebra system (in our case, Maple) to solve the system and find $r \leq \hat{\rho} := 0.6678$.

For the lower bound, first recall that
by Lemma \ref{lem:planeTreeCoeffIneq}     $B(z)$, the GF of full ordered binary trees, dominates coefficient-wise $T(z)$. Therefore the unique solution analytic at $0$ to the functional equation $y(z) = H(z, y(z))$, where  $H(z,w)$ is defined by: 
\begin{equation}\label{eq:upperBoundingSystem}
%{\small
\begin{split}
    H(z,w) &=  \frac{z}{2} 
        \exp\left( w + \sum\limits_{j = 2}^{N} \frac{\nicetrees(z^j)}{j} 
        + \sum\limits_{j > N} \frac{B(z^j)}{j} \right)  \\ 
      &+\frac{z}{2}\exp\left(-w + \sum\limits_{j = 2}^{N} (-1)^j \frac{\nicetrees(z^j)}{j} 
      + \sum\limits_{j > N} \frac{(-1)^j B(z^j)}{j} \right) \\&
       -\frac{z}{2}\left(w^2 + T(z^2) \right)
\end{split}
%}%
\end{equation}
dominates coefficient-wise $\nicetrees(z)$ and therefore its radius of convergence is at most $\rho$. Indeed, just recall that unrolling the exponentials, the negative signs are cancelled.  Also observe that 
\begin{equation}\label{eq:upperBoundingSystem2}
%{\small
\begin{split}
    H(z,w) &=  \frac{z}{2} 
        \exp\left( w + \sum\limits_{j = 2}^{N} \frac{L(z^j) +(T(z^j) - L(z^j))}{j}
        + \sum\limits_{j > N} \frac{B(z^j)}{j} \right)  \\ 
      &+\frac{z}{2}\exp\left(-w + \sum\limits_{j = 2}^{N} (-1)^j \frac{L(z^j) +(T(z^j) - L(z^j))}{j} 
      + \sum\limits_{j > N} \frac{(-1)^j B(z^j)}{j} \right) \\
      &-\frac{z}{2}\left(w^2 + L(z^2) +(T(z^2) - L(z^2))\right)
\end{split}
%}%
\end{equation}
 
Now, besides $L(z)$, the truncation of $T(z)$,  we also consider $M(z)$, the truncation of $B(z)$ up to powers at most $N$. Recall that $[z^{2n+1}]B(z)$ is the $n$-th Catalan number. 
We then  bound from above  
$\nicetrees(z^j) - L(z^j), j\geq 2,$  as well as $B(z^j),j \geq 2,$
 on the  interval $[0,\rho)$,  where all  $\nicetrees(z^j), j\geq 2,$ and $B(z^2)$ are analytic (because $\rho \leq  \hat{\rho} =  0.6678 < \sqrt(1/2)$), as follows.  

Observe that for $x \in [0,\rho)$ and $j\geq 2$, and because $\rho <1$ and  $B(z)$ bounds coefficient-wise $T(z)$, we have:

\begin{equation}\label{eq:error1}
	\nicetrees(x^j) - L(x^j) \leq B(x^j) - M(x^j)  \leq B(x^2) - M(x^2) \leq B(\hat{\rho}^2) -M(\hat{\rho}^2) :=e.
	\end{equation}

 By Lemma \ref{lemma:powerbounds},  we also have:

\begin{equation}\label{eq:error2}
     \frac{B(x^j)}{j} \leq  \frac{b^{j/2}}{j}
   \end{equation}
for $x \in [0,\hat{\rho})$, where $b := B(\hat{\rho}^2)$. 


We plug the errors given in  Equations \eqref{eq:error1} and \eqref{eq:error2} into the rhs of Equation \eqref{eq:upperBoundingSystem2}, to obtain the function:
\begin{equation}\label{eq:Gcheck}
\begin{split}
    \LG(z,w) &=  \frac{z}{2} 
        \exp\left( w + \sum\limits_{j = 2}^{N} \frac{L(z^j) +e}{j}
        + \sum\limits_{j > N} \frac{b^{j/2}}{j} \right)  \\ 
      &+\frac{z}{2}\exp\left(-w + \sum\limits_{j = 2}^{N} (-1)^j \frac{L(z^j) +e}{j} 
      + \sum\limits_{j > N} \frac{(-1)^j b^{j/2}}{j} \right) \\
      &-\frac{z}{2}\left(w^2 + L(z^2) +e\right).\\  
     % &=\frac{z}{2}M\exp\left( w + \sum_{j = 2}^{N} \frac{L(z^j)}{j}
       %  \right)  
           %  +\frac{z}{2}M\exp\left(-w + \sum_{j = 2}^{N} (-1)^j \frac{L(z^j)}{j} 
       %\right) \\
    %&-\frac{z}{2}\left(w^2 + L(z^2) +e\right), 
    \end{split}
 \end{equation} 
 Analogously to $\hat{G}$ we examined for the upper bound, we can show that the  entire $\check{G}$  satisfies  the conditions of Lemma \ref{lemma:implicitsing}. So let $\check{y}(z)$ be the analytic around 0 function implicitly defined by  $y(z) = \hat{G}(z, y(z)$ and also let $r,s = \lim_{x\rightarrow r^{-}}\check{y}(x)$ be the unique solution of the characteristic system $$\check{G}(r,s)=s, \ \check{G}_w(r,s) =1$$. 
 
 
 
 Now, if $h_{nm}$ and $\check{g}_{nm}$  are the Taylor coefficients of the two-variable functions 
  $H$ and   $\check{G}$, respectively, 
 we have that $\check{g}_{nm} \geq   h_{nm} $. Indeed, the appearance of  the error $e$ in $-\frac{z}{2}\left(w^2 + L(z^2) +e\right)$ vanishes when we differentiate. Therefore we may assume that the error $e$ appears in the exponentials only. So does $b$,  so $e$ and $b$  contribute multiplicatively to the coeffcients of $\check{G}$, so  $$\frac{\partial^n \partial^m \check{G}(0,0)}{\partial z^n \partial w^m} \geq \frac{\partial^n \partial^m H(0,0)}{\partial z^n \partial w^m}, \forall n, m.   $$
   
Therefore the coefficients $[z^j]\check{g}(z)$, which are given by a recurrence that entails $\check{g}_{nm}$, 
 are at least as large as the corresponding coefficients of the solution of $w = H(z,w)$, which are given by a recurrence that entails $h_{nm}$, 
 and therefore at least as large as the coefficients of $T(z)$, so $r \leq \rho$.
 
 
 
   
   We finally take, for example, $N=100$ and use Maple to find $r \geq \check{\rho} := 0.6677.$ \end{proof}   

\end{document}
